\documentclass[12pt]{ctexart}
\usepackage{amsmath, amsthm, amssymb, bm}
\usepackage[indent=0pt]{parskip}

\newcounter{probid}
\newcounter{subprobid}[probid]

\author{杨濛 2021213247}
\title{数学作业\\第三周}
\date{\today}

\newenvironment{problem}{
    \stepcounter{probid}
    \noindent\textbf{\arabic{probid}. }
}{\par}
\newenvironment{subproblem}{
    \stepcounter{subprobid}
    (\arabic{subprobid})
}{\par}
\newenvironment{solution}{
    \noindent\textbf{解: }
}{\par}

\begin{document}
    \maketitle
    \begin{problem}
        求下列幂级数的收敛域：
    \end{problem}
    \begin{subproblem}
        \(\sum\limits_{n=1}^\infty n^2 x^n\)
    \end{subproblem}
    \begin{solution}
        由于\(\lim\limits_{n\to\infty} \dfrac{n^2}{(n+1)^2} = 1\),所以原级数的收敛半径为1.

        当x=1时,原级数等于$\sum n^2$显然发散.故原级数的收敛域为$\left( -1, 1 \right)$.
    \end{solution}
    \begin{subproblem}
        \(\sum\limits_{n=1}^\infty \dfrac{2^n}{n^2+1} x^n\)
    \end{subproblem}
    \begin{subproblem}
        令\(a_n = \dfrac{2^n}{n^2+1}\), 则:
        \[\lim\limits_{n\to\infty} \dfrac{a_{n}}{a_{n+1}}
        = \dfrac{\left(n+1\right)^2+1}{2(n^2+1)} = \dfrac{1}{2}\]
        所以原级数的收敛半径为$\dfrac{1}{2}$.

        当$x=\dfrac{1}{2}$时,原级数等于$\sum\limits_{n=1}^\infty \dfrac{1}{n^2+1}$,收敛.

        当$x=-\dfrac{1}{2}$时,原级数绝对收敛.

        故原级数的收敛域为\(\left[-\dfrac{1}{2},\dfrac{1}{2}\right]\).
    \end{subproblem}
    \begin{subproblem}
        \(\sum\limits_{n=1}^\infty \left(-1\right)^n
        \dfrac{x^{2n+1}}{2n+1}\)
    \end{subproblem}
    \begin{solution}
        令\(a_n = \dfrac{\left(-1\right)^{n}}{2n+1}\),则 
        \(\lim\limits_{n\to\infty} \left|\dfrac{a_n}{a_{n+1}}\right| = 1\)

        故原级数收敛半径为1.
        
        当$|x|=1$时,原交错级数各项绝对值单调递减,是收敛的莱布尼茨级数.

        综上,该级数收敛域为$\left[-1, 1\right]$.
    \end{solution}
    \begin{subproblem}
        \(\sum\limits_{n=1}^\infty n!\left(x-1\right)^n\)
    \end{subproblem}
    \begin{solution}
        令$a_n = n!\left(x-1\right)^n$,则
        \[\lim\limits_{n\to\infty} \dfrac{a_{n+1}}{a_n}
        = \lim\limits_{n\to\infty } (n+1)\left(1-\dfrac{1}{n}\right)^n
        = \lim\limits_{n\to\infty} \dfrac{n}{\rm e} = +\infty\]
        所以原级数的收敛半径为0, 收敛域为$\{0\}$.
    \end{solution}
    \begin{subproblem}
        \(\sum\limits_{n=1}^\infty 
        \dfrac{3^n}{n!}
        \left(\dfrac{x-1}{2}\right)^n\)
    \end{subproblem}
    \begin{solution}
        令\(a_n = \dfrac{3^n}{n!}\),则
        \[\lim\limits_{n\to\infty} \dfrac{a_{n+1}}{a_n}
        = \lim\limits_{n\to\infty} \dfrac{3}{n+1} = 0\]
        所以原级数在$\dfrac{x-1}{2}\in\mathbf R
        \Rightarrow x \in \mathbf{R}$上收敛.
    \end{solution}
    \begin{subproblem}
        \(\sum_{n=1}^\infty \left[
            \dfrac{\left(-1\right)^n}{n}
            + \dfrac{1}{2^n}
        \right] x^n\)
    \end{subproblem}
    \begin{solution}
        转化为两个级数\(\sum\limits_{n=1} \dfrac{(-1)^n}{n}x^n\)
        和\(\sum\limits_{n=1} \dfrac{1}{2^n} x^n\).

        对于\(\sum\limits_{n=1} \dfrac{(-1)^n}{n}x^n\), 
        由于\[\lim\limits_{n\to\infty} \left|-\dfrac{n}{n-1}\right| = 1\],
        故该级数收敛半径为1.

        当$x=1$时,该部分为收敛的莱布尼茨级数;当$x=-1$时,该部分为发散的调和级数.

        对于\(\sum\limits_{n=1} \dfrac{1}{2^n} x^n\), 
        易知当$|x<2|$时该部分收敛.

        由级数性质可知,原级数在$\left(-1, 1\right]$上收敛.
    \end{solution}
    \begin{problem}
        求下列幂级数的和函数:
    \end{problem}
    \begin{subproblem}
        $\sum\limits_{n=1}^\infty nx^n$
    \end{subproblem}
    \begin{solution}
        该级数的收敛域为$(-1,1)$.令和函数为$S(x)$,则
        \[S(x) = \sum\limits_{n=1}^\infty nx^n
        = x\sum\limits_{n=1}^\infty nx^{n-1}\]
        令$T(x) = \sum\limits_{n=1}^\infty nx^{n-1}$, 则
        \[\int^x_0 T(x) = \sum_{n=1}^\infty x^n {\rm d}x
        = \dfrac{1}{1-x}\]
        所以有
        \[T(x) = \dfrac{1}{(1-x)^2}{\rm d}x
        \Rightarrow S(x) = \dfrac{x}{(1-x)^2}\]
    \end{solution}
    \begin{subproblem}
        \(\sum\limits_{n=1}^\infty \dfrac{n(n+1)}{2} x^{n-1}\)
    \end{subproblem}
    \begin{solution}
        易得该级数的收敛域为(-1, 1).令该级数的和函数为$S(x)$.

        注意到\(\dfrac{n(n+1)}{2} x^{n-1} {\rm d}^2x
        = \dfrac{{\rm d}^2}{{\rm d}x^2} \dfrac{x^{n+1}}{2}\),所以
        \[S(x) = \dfrac{{\rm d}^2}{{\rm d}x^2}
        \sum\limits_{n=1}^\infty \dfrac{x^{n+1}}{2}
        = \dfrac{{\rm d}^2}{{\rm d}x^2}
        \dfrac{x^2}{1-x}
        = \dfrac{x}{(1-x)^2}\]
    \end{solution}
    \begin{problem}
        求级数\(\sum\limits_{n=1}^\infty \dfrac{2n+1}{3^n}\)的和.
    \end{problem}
    \begin{solution}
        令\(S(x) = \sum\limits_{n=1}^\infty (2n+1)x^n\),
        则原级数的收敛半径为$1$.

        注意到
        \[S(x) = \sum\limits_{n=1}^\infty [2(n+1)-1]x^n
        = \sum\limits_{n=1}^\infty 2(n+1)x^n
        - \sum\limits_{n=1}^\infty x^n\]
        而
        \begin{align*}
            \int_0^x \sum\limits_{n=1}^\infty 2(n+1)x^n
            & = \int_0^x \sum\limits_{n=1}^\infty 2(n+1)x^n \\
            & = \sum\limits_{n=1}^\infty 2x^{n+1} \\
            & = \dfrac{2x^2}{1-x}
        \end{align*}
        所以
        \begin{equation*}
            \sum\limits_{n=1}^\infty 2(n+1)x^n
             = \dfrac{4x}{1-x}
        \end{equation*}
        因此有
        \begin{align*}
            S(x) & = \sum\limits_{n=1}^\infty [2(n+1)-1]x^n \\
            & = \dfrac{4x}{1-x} - \dfrac{x}{1-x} \\
            & = \dfrac{3x}{1-x}
        \end{align*}
        当$x=\dfrac{1}{3}$时,级数等价于
        \(\sum\limits_{n=1}^\infty \dfrac{2n+1}{3^n}
        = \dfrac{3}{2}\).
    \end{solution}
    \begin{problem}
        将下列函数展开为x的幂级数,并写出展开式成立的区间.
    \end{problem}
    \begin{subproblem}
        \(\ln\left(a+x\right)\)
    \end{subproblem}
    \begin{solution}
        \[\ln(a+x) = \ln\left(1+\dfrac{x}{a}\right) + \ln a\]
        在$\left|\dfrac{x}{a}\right| < 1 
        \Rightarrow \left| x \right| < a$的范围内, 上式可以展开为:
        \[\ln a + \sum\limits_{n=1}^\infty 
        (-1)^{n-1}\dfrac{x^n}{na^n}\]
    \end{solution}
    \begin{subproblem}
        $\sin^2 x$
    \end{subproblem}
    \begin{solution}
        \begin{align*}
            \sin^2 x & = \dfrac{1-\cos 2x}{2}\\
            & = \dfrac{1}{2} \sum\limits_{n=1}^\infty
            (-1)^{n-1} \dfrac{x^{2m}}{(2n)!}
        \end{align*}
        对于\(x\in\mathbf{R}\)均成立.
    \end{solution}
    \begin{subproblem}
        $\left(4-x^2\right)^{-\frac{1}{2}}$
    \end{subproblem}
    \begin{solution}
        \begin{align*}
            & \left(4-x^2\right)^{-\frac{1}{2}}\\
            = & \dfrac{1}{2} \left(
                1-\dfrac{x^2}{4}
            \right)\\
            = & \dfrac{1}{2} \left( 1 + \sum\limits_{n=1}^\infty
            \dfrac{m(m-1)\cdots(m-n+1)}{n!} 
            \left(-\dfrac{x^2}{4}\right)^n\right)
        \end{align*}
        当$x\in (-1,1)$时成立.
    \end{solution}
    \begin{subproblem}
        \(\dfrac{x}{\sqrt{1+x^2}}\)
    \end{subproblem}
    \begin{solution}
        注意到\(\dfrac{x}{\sqrt{1+x^2}}
        = \left(\sqrt{1+x^2}\right)^\prime\), 而
        \begin{align*}
            & \sqrt{1+x^2}\\
            = & 1 + \dfrac{x^2}{2} - \dfrac{1}{2^2\times 2!}x^4 
            + \dfrac{3}{2^3\times 3!}x^6 + \cdots\\
            = & 1 + \dfrac{1}{2}x^2 + \sum\limits_{n=1}^\infty
            (-1)^{n-1} \dfrac{(2n-1)!!}{(2n+2)!!}x^{2n+2}\\
        \end{align*}
        所以
        \begin{align*}
            \sqrt{1+x^2} = 
             x + \dfrac{1}{6}x^3 + \sum\limits_{n=1}^\infty
            (-1)^{n-1} \dfrac{(2n-1)!!}{(2n)!!} x^{2n+3}
        \end{align*}
        由课本例题可知该级数在$x\in [-1,1]$处收敛.
    \end{solution}
    \begin{subproblem}
        \(\dfrac{x}{(1-x)(1-2x)}\)
    \end{subproblem}
    \begin{solution}
        
    \end{solution}
\end{document}